\documentclass{book}
\usepackage[utf8]{inputenc}
\usepackage[T1]{fontenc}
\usepackage{geometry}
\usepackage{graphicx}
\usepackage{float}
\usepackage{hyperref}
\usepackage[english,indonesian]{babel}

\title{Optimasi Pemrosesan Template LaTeX Menggunakan Pendekatan Berbasis AI}
\author{Academic Support Agent}
\date{\today}

\begin{document}
\maketitle

\begin{abstract}

\end{abstract}

\chapter{PENDAHULUAN}


\chapter{TINJAUAN PUSTAKA}


% Original content (commented out):
% \documentclass[11pt,a4paper, oneside]{book}
\usepackage[fontsize=11.05pt]{fontsize}
\usepackage{thesis_espoch}
\usepackage[utf8]{inputenc}
\usepackage[spanish, english]{babel}
\usepackage{geometry}
\usepackage{amsmath}
\usepackage{amsfonts}
\usepackage{amssymb}
\usepackage{graphicx}
\usepackage{setspace}
\spacing{1.5}
\usepackage[hidelinks]{hyperref}
\usepackage{wrapfig}
\usepackage[rightcaption]{sidecap}
\usepackage{times}
\usepackage{physics}
\usepackage{xcolor,colortbl}
\usepackage{multirow}
\usepackage{ragged2e}
\usepackage{threeparttable}
\usepackage[mathscr]{euscript}
\usepackage{subfigure} 
\setcounter{tocdepth}{3}
\usepackage[pagestyles]{titlesec}
\usepackage{titletoc}
\usepackage{csquotes}
\usepackage[style=iso-authoryear,currentlang,backend=bibtex]{biblatex}

\addbibresource{Referencias1.bib} 

% LEER CON DETENIMIENTO LOS SIGUIENTES COMENTARIOS
%#################################################################
%#################################################################

% ¿Cómo agregar archivo de bibliografías?
%\addbibresource{Referencias1.bib} 

% ¿CÓMO CITAR BIBLIOGRAFÍAS?
% \textcite ---> Para citar como texto, ejemplo: Según XXXX, 20XX.
% \parencite  --> Para citar con paréntesis, ejemplo: (XXXX, 20XX)

%#################################################################
%#################################################################

% POSIBLES ERRORES 
% ¿CÓMO HACER BREAKPAGE EN CASO DE QUE FALLE LOS INDICES DE ILUSTRACIONES O TABLA?
%\addtocontents{lof}{\protect\pagebreak} -> agregar esta linea de comando en después de la imagen en la cual se quiere hacer el salto de página, esto es para figuras. Para tablas sería \addtocontents{lot}{\protect\pagebreak}

%#################################################################
%#################################################################

% DATOS DE LAS PORTADAS___________________________________________
% Respete las mayúsculas y minúsculas de los ejemplos
\Facultad{Facultad de Ciencias}
\Carrera{Carrera de Física}
% Tema de tesis
\title{Xxxxxx Xxxxxx Xxxxxx Xxxxxx Xxxxxx Xxxxxx Xxxxxx Xxxxxx}
% Nombre del autor/a pero respetándo las maúsculas y minúscula.
\author{Xxxxxx Xxxxxx Xxxxxx Xxxxxx}
% Supervisor/a, nombre en mayúsculas
\supervisor{Dr. XXXXXX XXXXXX XXXXX}
% Tipo de proyecto
\project{Proyecto de investigación}
\titulo{FÍSICO}
% Año en curso
\date{2024}

% PARA EL RESTO DE HOJAS PRELIMINARES_____________________________

% Autor o autora en minúscula
\autoroautora{autora} 
% el señor o la señora en minúscula 
\arquetipo{el señor}
% Número de cédula 
\cedula{1234567890}  
% Fecha de defensa de la tesis
\DefensaDate{1 de enero de 2024} 
\FDDate{2024-02-01} %aaa-mm-dd

% MIEMBROS DEL TRIBUNAL (minúscula), respetar el cargo, solo cambiar género de cargo
\NameT{Fis. Juanito Alimaña PhD.}
\CargoT{PRESIDENTE DEL TRIBUNAL}

\NameTT{Fis. Juanito Alimaña PhD.}
\CargoTT{DIRECTOR DEL TRABAJO DE INTEGRACIÓN CURRICULAR}

\NameTTT{Fis. Juanito Alimaña PhD.}
\CargoTTT{ASESOR DEL TRABAJO DE INTEGRACIÓN CURRICULAR}

%#################################################################
%#################################################################

\usepackage{lipsum} % NO ES NECESARIO TENER ESTO

\begin{document}

\maketitle

\newpage
\thispagestyle{plain}
\begin{center}
\textbf{DEDICATORIA}
\vspace{1cm}
\end{center}
\lipsum[100] 

\newpage
\thispagestyle{plain}
\begin{center}
\textbf{AGRADECIMIENTO}
\vspace{1cm}
\end{center}
\lipsum[100] 
\newpage

% NO BORRAR/MODIFICAR ESTA SECCIÓN DE CÓDIGO AL MENOS QUE SEPA LO QUE ESTÁ HACIENDO.
\renewcommand\tablename{\textbf{Tabla}}
\renewcommand\figurename{\textbf{Ilustración}}
\renewcommand*\contentsname{\normalsize ÍNDICE DE CONTENIDO}
\renewcommand{\listfigurename}{\normalsize \hspace{-0.2cm}ÍNDICE DE ILUSTRACIONES}
\renewcommand{\listtablename}{\normalsize ÍNDICE DE TABLAS}
\tableofcontents
\newpage
\listoftables
\addcontentsline{toc}{section}{\textbf{ÍNDICE DE TABLAS}}
\newpage
\listoffigures
\addcontentsline{toc}{section}{\textbf{ÍNDICE DE ILUSTRACIONES}}
\newpage
\listofappendices
\addcontentsline{toc}{section}{\textbf{ÍNDICE DE ANEXOS}}

% RESUMENES EN ESPAÑOL E INGLÉS
%#################################################################
%#################################################################

\newpage
\pagestyle{plain}
\addcontentsline{toc}{section}{RESUMEN}
\begin{Resumen}
\lipsum[100] \lipsum[100] \lipsum[100] \\

\textbf{Palabras clave:} $<$PALABRA$>$, $<$PALABRA$>$, $<$PALABRA$>$, $<$PALABRA$>$.
\end{Resumen}

\newpage
\addcontentsline{toc}{section}{SUMMARY/ABSTRACT}
\pagestyle{plain}
\begin{Abstract}
\lipsum[100] \lipsum[100] \lipsum[100]\\

\textbf{Keywords:} $<$WORD$>$, $<$WORD$>$, $<$WORD$>$, $<$WORD$>$. 
\end{Abstract}
\newpage

%#################################################################
%#################################################################

\renewcommand{\thepage}{\arabic{page}}
\setcounter{page}{1}
\pagestyle{plain}

\section*{INTRODUCCIÓN}
\addcontentsline{toc}{section}{\textbf{INTRODUCCIÓN}}

\lipsum[6,paragraph]\\

\lipsum[6,paragraph]\\

\lipsum[6,paragraph] \parencite{C1}.

\part*{CAPÍTULO I} 
\addcontentsline{toc}{part}{CAPÍTULO I}

\chapter{MARCO TEÓRICO PREFERENCIAL}

\section{Antecedentes}

\subsection{Xxxxx Xxxxx Xxxxx}

\lipsum[6,paragraph]\\

\lipsum[6,paragraph]\\

\lipsum[6,paragraph]

% EJEMPLO DE IMAGEN
\begin{figure}
\centering\captionsetup{singlelinecheck=false,width=0.45\linewidth}  
\includegraphics[scale=0.2]{ejemplo.jpg}
\caption{XXXXX XXXXXX XXXXXXX}
\begin{flushleft}
\hspace{4.15cm}\scriptsize \textbf{Realizado por:} XXXXX, X. X., 202X. % procura siempre modificar el espacio de este texto, debe estar al marguen de la caption.
\end{flushleft}
\label{Imagen_1}
\end{figure}

% EJEMPLO TABLA
\begin{table}
\centering
\begin{threeparttable}
\caption{Resultado datos}
\label{Tabla_1}
\begin{tabular}{lcrrrr}
\hline
Alpha & Beta & Gamma & Theta & Delta \\
\hline\noalign{\smallskip}

\hline\noalign{\smallskip}
\end{tabular}
\raggedright\scriptsize \textbf{Realizado por:} XXXXX, XXXX, 202X.
\end{threeparttable}
\end{table}

\lipsum[6,paragraph]\\

\lipsum[6,paragraph]\\

\lipsum[6,paragraph]

\newpage
\chapter*{BIBLIOGRAFÍA}
\addtocontents{toc}{\vspace{0.75cm}\textbf{BIBLIOGRAFÍA}\\}
\pagestyle{empty}
\vspace{-0.5cm}
\defbibheading{myheading}[References]{}
\printbibliography[heading=myheading]


\newpage
\chapter*{ANEXOS}
\appendix 
\addtocontents{toc}{\textbf{ANEXOS}}
\vspace{-1cm}


\chapter*{ANEXO A: XXXXXXX XXXXXX XXXXXXXXXX}
\addtocontents{apc}{\textbf{ANEXO A:} XXXXXXX XXXXXX XXXXXXXXXX\\}  % Para agregar al índice de Anexos


\end{document}

% ¿CÓMO INGRESAR IMÁGENES?  ______________________________________________________________

\begin{figure}
\centering\captionsetup{singlelinecheck=false, width=0.80\linewidth}  % regula tamaño del box de caption
\includegraphics[scale=0.64]{CCF.png}
\caption{XXXX XXXXXX}
\begin{flushleft}
\hspace{2.4cm}\scriptsize \textbf{Realizado por:} XXXXX, X. X., 202X. % modificar el \hspace a gusto dependiendo de la imagen. 
\end{flushleft}
\label{Ima_4}
\end{figure}

% ¿CÓMO HACER TABLAS? ____________________________________________________________________
\begin{table}
\centering
\begin{threeparttable}  % permite la justificación a la izquierda del caption
\caption{XXXXX XXXXXX XXXXX}
\label{Tabla_1}
\begin{tabular}{lcrrrr}
\hline
Alpha & Beta & Gamma & Theta & Delta \\
\hline\noalign{\smallskip}

\hline\noalign{\smallskip}
\end{tabular}
\raggedright\scriptsize \textbf{Realizado por:} XXXXX, XXXX, 202X.
\end{threeparttable}
\end{table}

% ¿COMO HACER SUBFIGURAS?____________________________________________________________________
\begin{figure}
\centering\captionsetup{singlelinecheck=false,width=0.85\linewidth}
\subfigure[Subtexto de la subfigura]{\includegraphics[scale=0.15]{ejemplo1.png}\label{Ima_4}}\hspace{0.5cm}
\subfigure[Subtexto de la subfigura]{\includegraphics[scale=0.27]{ejemplo1.png}\label{Ima_5}}
\caption{}\label{Ima_6}
\begin{flushleft}
\hspace{1.15cm}\scriptsize \textbf{Realizado por:} XXXXX, XXXX, 202X.
\end{flushleft}
\end{figure}

% Para referenciar los subplots hacemos: 
\ref{Ima_6}\subref{Ima_5} o \ref{Ima_6}\subref{Ima_4}


\end{document}